\begin{lemma}[Lexicographic Step 1 exit-leg minimum]
Fix $G,k,e_0,n,q$.  Suppose that the family of Step 1 exit-leg
presentations of size $(n,q)$ is nonempty and that the terminal vertex
set of every member is elementary.  Then this family contains a Step 1
exit-leg minimum.
\end{lemma}
\begin{proof}
For a presentation $R$ in the given nonempty family, let $d(R)$ be the
divider count of its good hierarchy, one of the fields listed in
\noderef{Step1ExitLegPresentation}.  The nonempty set of natural numbers
$d(R)$ has a least member; call it $d_0$ and restrict attention to
presentations with $d(R)=d_0$.

Define a second natural-number measure $s(R)$ on this restricted family.
If the second color $R.b$ is missing on $R.T$, put $s(R)=0$.  Otherwise
apply \noderef{ElementaryClosedExitPairExtendsTriple}.  Its hypotheses
are supplied by the assumed elementariness of $R.T$ and by the
properness, distinct colors, missing and singleton-closed first color,
and two stored support-disjoint exit legs in
\noderef{Step1ExitLegPresentation}.  It gives a third exit leg
support-disjoint from the stored two, with the asserted endpoint
inequalities and with an exact $R.b$-colored boundary edge on that leg.
Thus the set of totals
\[
 R.P_1.\mathit{length}+R.P_2.\mathit{length}+P_3.\mathit{length}
\]
over all such third-leg packages is nonempty.  Let $s(R)$ be one plus
the least such total.  The set of $s(R)$ over all $d_0$-presentations is
again a nonempty set of natural numbers.  Choose $P$ for which $s(P)$ is
least.

Suppose first that $s(P)>0$.  Then $P.b$ is not missing on $P.T$.
Choose a third package $(a_3,b_3,P_3,f_3)$ attaining the least total in
the definition of $s(P)$.  For the first two packages use the exact
stored edges $P.f_1,P.f_2$ from
\noderef{Step1ExitLegPresentation}; for the third use the exact edge
$f_3$ furnished by \noderef{ElementaryClosedExitPairExtendsTriple}.
Consequently, for each $i=1,2,3$, the complete quadruple
$(a_i,b_i,P_i,f_i)$ contains an exit path, its two ends, and an exact
$P.b$-colored boundary edge lying on that path and incident with its
tree end.

The three tree ends are distinct, because the three path supports are
pairwise disjoint and each path contains its tree-end vertex-node.
By \noderef{TreeExitPath}, every $b_i$ belongs to $V(P.T)$; by
\noderef{TreeSequenceVertexSet} this is the set of displayed vertices.
Those vertices occur at unique positions because the vertex list in
\noderef{TreeSequence} is injective.  Sort the three complete quadruples,
not merely the paths, by these positions.  If the resulting order is
$(a'_1,b'_1,P'_1,f'_1),(a'_2,b'_2,P'_2,f'_2),
(a'_3,b'_3,P'_3,f'_3)$, then
\noderef{TreeVertexOrder} gives $b'_1\prec b'_2\prec b'_3$.
Because the permutation transports the whole quadruple, every $f'_i$
is still the boundary edge on $P'_i$ incident with $b'_i$.

Construct a presentation $\widehat P$ by retaining all data of $P$
through its colors and replacing its first two exit-leg packages by the
first two sorted quadruples.  Pairwise support-disjointness gives the
required disjointness and endpoint inequalities, and the preceding
paragraph gives the strict tree-end order and both retained boundary
identities.  Thus every exit-leg field of
\noderef{Step1ExitLegPresentation} holds.  All earlier fields are copied
from $P$, so $d(\widehat P)=d_0$.  Its sorted third quadruple is a
qualifying third package of the same total length as the original
triple.  Hence $s(\widehat P)\leq s(P)$.  Global minimality of $s(P)$
gives the reverse inequality, so equality holds.  In particular, if the
relabeled first two quadruples admitted a qualifying third leg of
smaller total length, they would form this same comparison presentation
$\widehat P$ with $s(\widehat P)<s(P)$, a contradiction.  Therefore the
sorted triple is least for $\widehat P$.  Replace $P$ by $\widehat P$;
we have obtained the ordered three-leg alternative required below,
with its correct transported edge $f'_3$.

We now establish primary minimality.  Take arbitrary data occurring in
the primary comparison of \noderef{Step1ExitLegMinimum} and suppose its
two exit paths exist.  By
\noderef{CorrespondingExtendedTashkinovTree}, the comparison carries an
actual replacement input, an actual ETT, and the same rung number and
certified series as the selected ETT.  The selected rung equality in
\noderef{Step1ExitLegPresentation} therefore makes the replacement rung
number $n$.  The comparison hypotheses supply maximum property,
properness, the exact-one-uncolored condition, its large bound, and its
terminal minus-closed good hierarchy.  Its divider count is less than
$d_0$, hence also less than $q$.

Apply \noderef{ClosedColorExitPathBoundaryEdge} separately to the two
comparison exit paths, using their common singleton-closed first color.
This supplies a complete quadruple, including the exact second-color
boundary identity, for each path.  Their tree ends are distinct because
a common tree end would be a vertex-node in both disjoint supports.
As above, \noderef{TreeExitPath}, \noderef{TreeSequenceVertexSet}, and
the injectivity in \noderef{TreeSequence} give two distinct displayed
positions.  Sort the two complete quadruples by those positions;
\noderef{TreeVertexOrder} gives the required strict order while the
boundary identities remain attached to their own paths.  Together with
the comparison data, these are exactly all fields of
\noderef{Step1ExitLegPresentation}.  They therefore define a
presentation whose divider count is less than $d_0$, contradicting the
definition of $d_0$.  This proves the primary certificate.

For secondary minimality, take arbitrary data in the secondary
comparison and normalize its first two paths in exactly the same way:
apply \noderef{ClosedColorExitPathBoundaryEdge} to both paths and sort
the two complete quadruples.  The replacement ETT and all non-leg fields
are justified exactly as in the primary comparison; equality of divider
counts makes the resulting
\noderef{Step1ExitLegPresentation} $R$ a $d_0$-presentation.  If its
second color is missing, then $s(R)=0<s(P)$, since this certificate is
needed only in the positive case.  Otherwise the comparison supplies a
third path support-disjoint from the two and with distinct outside end.
Support-disjointness also makes its tree end distinct; applying
\noderef{ClosedColorExitPathBoundaryEdge} to it supplies its exact
second-color boundary identity.  Thus it is one of the qualifying third
packages used to define $s(R)$.  Sorting the first two packages did not
change their lengths, so if the displayed total is smaller than the
selected sorted total, then
\[
 s(R)\leq 1+\text{(displayed total)}<s(P),
\]
again contradicting the choice of $P$.  This proves the secondary
certificate.

If the originally selected secondary value was zero, keep that selected
presentation; its second color is missing and the secondary certificate
is the first alternative.  If it was positive, use the sorted
presentation constructed above and its ordered least triple.  In either
case attach the elementariness assumed for every presentation and the
primary and secondary certificates just proved.  The definition
\noderef{Step1ExitLegMinimum} then gives a Step 1 exit-leg minimum.
\end{proof}
